% ============================================================================
%  基础物理实验 实验报告 模板（LaTeX / XeLaTeX）
%  格式取自课程推荐模板（UCAS 基础物理实验 Word 模板的 LaTeX 版）
%
%  编译：  xelatex 实验报告.tex       跑两遍，交叉引用与页码才准
%          （或 latexmk -xelatex 实验报告.tex）
%
%  用法：把本文件复制到 <实验>/实验数据/实验报告.tex，改下面「报告头」和正文。
% ============================================================================
\documentclass[11pt]{article}

\usepackage[a4paper]{geometry}
\geometry{left=2.0cm,right=2.0cm,top=2.5cm,bottom=2.5cm}

\usepackage{ctex}                       % 中文支持
\usepackage{amsmath,amssymb,bm,mathrsfs,mathtools,breqn}
\usepackage{graphicx,subfigure}
\usepackage{booktabs,multirow,bigstrut,array}
\usepackage{float}                      % 提供 [H] 浮动体位置
\usepackage{fancyhdr}
\usepackage{tikz,xcolor}
\usepackage{pdfpages}                   % 插原始记录表扫描件
\usepackage[capitalize]{cleveref}       % \cref 自动加前缀
\Crefname{section}{Section}{Sections}
\Crefname{table}{Table}{Tables}
\Crefname{figure}{Figure}{Figures}

% ---- 字体 ------------------------------------------------------------------
% 默认用系统字体（Windows / 装了中文字体的机器都能直接编）。
% 换机器若报缺字体，把课程模板文件夹里的 ttf 拷到本目录，改成下面两行：
%   \setmainfont{Palatino Linotype.ttf}
%   \setCJKmainfont{SimHei.ttf}
\setmainfont{Palatino Linotype}
\setCJKmainfont{SimHei}
\punctstyle{kaiming}
\renewcommand{\emph}[1]{\begin{kaishu}#1\end{kaishu}}

% ============================================================================
%  报告头：每个实验改这一段
% ============================================================================
\newcommand{\experiName}{测量金属的杨氏模量}
\newcommand{\supervisor}{XXX}
\newcommand{\name}{XXX}
\newcommand{\studentNum}{20XXXXXXXXXXX}
\newcommand{\class}{1}                  % 分班
\newcommand{\group}{04}                 % 分组
\newcommand{\seat}{05}                  % 座号
\newcommand{\dateYear}{2026}
\newcommand{\dateMonth}{9}
\newcommand{\dateDay}{21}
\newcommand{\room}{XXX}
\newcommand{\others}{$\square$}         % 调课/补课：是则改 $\square$\hspace{-1em}$\surd$
% ============================================================================

\begin{document}

\begin{center}
    \LARGE \bf 《\, 基\, 础\, 物\, 理\, 实\, 验\, 》\, 实\, 验\, 报\, 告
\end{center}

\begin{center}
    \noindent \emph{实验名称}\underline{\makebox[25em][c]{\experiName}}
    \emph{指导教师}\underline{\makebox[8em][c]{\supervisor}}\\
    \emph{姓名}\underline{\makebox[6em][c]{\name}}
    \emph{学号}\underline{\makebox[10em][c]{\studentNum}}
    \emph{分班分组及座号} \underline{\makebox[5em][c]{\class \ -\ \group \ -\ \seat}\emph{号}}\\
    \emph{实验日期} \underline{\makebox[3em][c]{\dateYear}}\emph{年}
    \underline{\makebox[2em][c]{\dateMonth}}\emph{月}
    \underline{\makebox[2em][c]{\dateDay}}\emph{日}
    \emph{实验地点}\underline{\makebox[4em][c]{\room}}
    \emph{调课/补课} \underline{\makebox[3em][c]{\others\ 是}}
    \emph{成绩评定} \underline{\hspace{5em}}
    {\noindent}
    \rule[8pt]{17cm}{0.2em}
\end{center}

% ============================================================================
%  正文。章节顺序按课程模板；实验没有的节（如思考题）整节删掉即可。
%  文风底线：不列中间代入算式、括号不塞废话、不复述表格里已有的数字。
% ============================================================================

\section{实验目的}
\begin{enumerate}
    \item ...
\end{enumerate}

\section{实验原理}
% 公式用 equation（带编号）或 \[ \]（不带）。多行对齐用 align。
\begin{equation}
    Y = \frac{4MgL}{\pi d^{2}\Delta L}
\end{equation}

\section{实验仪器与用具}
% 写清型号；指导书写了技术指标的可照抄要点。

\section{实验步骤}
\begin{enumerate}
    \item ...
\end{enumerate}

\section{实验结果与数据处理}

% --- 表格规范（数据表、曲线图要素示意）--------------------------------------
% 三线表。caption 要含「表号 + 测量条件」，如「表一：摆长为 70\,cm 时…」；
% 表头写「物理量符号 / 单位」；表里只放实验数据，不放算出来的中间量。
\begin{table}[H]
    \centering
    \caption{<测量条件>下 <物理量> 的测量结果}
    \begin{tabular}{ccccc}
        \toprule
        $i$ & $M/\mathrm{g}$ & $l/\mathrm{mm}$ & $l'/\mathrm{mm}$ & $\bar l/\mathrm{mm}$ \\
        \midrule
        1 & 250 & $-0.10$ & $-0.05$ & $-0.075$ \\
        \bottomrule
    \end{tabular}
\end{table}

% --- 图规范 ----------------------------------------------------------------
% 图里必须有：图例、数据点＋拟合曲线、拟合表达式与 $R^2$（画在图上）、
% 带单位的坐标轴标签、合适标度；caption 写「物理量随物理量的变化关系」。
% 拟合式与 $R^2$ 由处理脚本画进 PNG，不在正文里再抄一遍。
\begin{figure}[H]
    \centering
    \includegraphics[width=0.7\textwidth]{fig1.png}
    \caption{叉丝偏移量 $\delta L$ 随砝码质量 $M$ 的变化关系}
\end{figure}

\section{思考题}
% 指导书【思考题】原题抄一条、答一条；没有思考题的实验整节删掉。
\begin{enumerate}
    \item <原题>

    答：<回答>
\end{enumerate}

\section{总结}

\section{实验原始数据记录表}
% 记录纸拍照 → 转 PDF 后整页插入。多张照片先合成一个 pdf。
\includepdf[pages=-]{记录表.pdf}

\end{document}
